%------------------------------------------------------------------------------%
\documentclass[fleqn,utf8,aspectratio=169,12pt]{beamer}

\usepackage{calculo_varias_variaveis-1}

\title{Vetores}

%------------------------------------------------------------------------------%
\begin{document}

\addtitle

%------------------------------------------------------------------------------%
\section{Vetores}

%------------------------------------------------------------------------------%
\begin{frame}{Vetor}
\begin{itemize}[<+->]
  \setlength\itemsep{1em}
  \item Vetor é definido por \emph{módulo}, \emph{direção} e \emph{sentido}
  \item Representa uma classe de \emph{segmentos de reta orientados} equipolentes\\
        (paralelos, mesmo tamanho e sentido)
  \item É representado por componentes
        \(\vetor{v_1}{v_2}\) ou \(\vetortri{v_1}{v_2}{v_3}\)
  \item Pode ser posicionado em qualquer ponto \((x, y)\) ou \((x, y, z)\)
\end{itemize}
\end{frame}

\include{ex/B-03-Vetores-ex1}

%------------------------------------------------------------------------------%
%% \section{Lista Mínima}
%%
%% %------------------------------------------------------------------------------%
%% \begin{frame}{Lista Mínima}
%%   Cálculo Vol. 2 do Thomas 12\textsuperscript{a} ed. -- Seção 11.2
%%   \begin{enumerate}
%%     \item Estudar todo o texto da seção
%%     \item \emph{Calcule o vetor tangente}
%%           das funções dos exercícios: 2, 8, 12, 16, 18, 20\\
%%           (Ignore o enunciado original dos exercícios)
%%   \end{enumerate}
%%
%%   \vfill
%%   Atenção: A prova é baseada no livro, não nas apresentações
%% \end{frame}

%------------------------------------------------------------------------------%
\end{document}
%------------------------------------------------------------------------------%
